% FILE: 00_project_identity.tex
% PROJECT: sources-papers
% ROLE: academic-paper-classification-and-evidence-grounding

\section{Project Identity}
\label{sec:sources-papers-identity}

\subsection{Purpose}
\label{sec:sources-papers-identity-purpose}

The \texttt{sources/papers} directory is the academic grounding layer of the
process-intelligence research foundry. Its governing purpose is to translate
formal results published in the process mining and workflow literature into
compile-time obligations --- Rust type laws, execution graduation boundaries,
and board-admissible evidence claims --- that downstream components
(\texttt{wasm4pm-compat} and \texttt{wasm4pm}) must satisfy before any
certification claim can be issued.

The corpus does not merely catalogue papers. Each paper is assessed along
three independent obligation axes: (1) what formal objects it introduces that
must exist as zero-cost Rust types in \texttt{wasm4pm-compat}; (2) what
algorithms it mandates that must execute in the \texttt{wasm4pm} engine with
typed input/output boundaries; and (3) what board-admissible claims it enables
when both type obligations and execution obligations are met. A paper for which
axis~(1) is satisfied but axis~(2) is absent carries the status
\texttt{COVERED\_BY\_GRADUATION\_BOUNDARY}, not \texttt{COVERED\_BY\_TYPE}.
The distinction is precise and consequential: no claim graduates to
board-admissible until the execution boundary is crossed.

This layer was the authority source for all downstream refactors of
\texttt{wasm4pm} and \texttt{wasm4pm-compat}. The commitment documented in
\texttt{CLAUDE.md} is explicit: no downstream \texttt{wasm4pm} refactor may
proceed until the research program speaks; no M\&A claim may be manufactured
without research program grounding.

\subsection{Repository Surface}
\label{sec:sources-papers-identity-surface}

The paper corpus consists of 15 files at assessment date (2026-05-31), against
a gate target of 7 --- the surplus representing depth rather than mere
compliance. The primary inventory document is
\texttt{paper-canon.md} (version 30.1.2, status COMPLETE), which provides
conformance claims for 9 papers with 150+ formal objects, 75+ algorithms
catalogued, 120+ failure conditions, and 25+ board claim archetypes.

Three mapping documents partition the obligation surface:
\texttt{PAPER\_TO\_TYPE\_LAW.md} enumerates the zero-cost Rust type surface
required by each paper; \texttt{PAPER\_TO\_EXECUTION\_LAW.md} enumerates
algorithm graduation requirements with input/output type shapes;
\texttt{PAPER\_TO\_BOARD\_CLAIM.md} identifies which claims are board-admissible
now versus which are blocked by missing types or missing execution coverage.

Seven individual deep-analysis files cover the most consequential algorithms:
Alpha Miner, Inductive Miner, OCEL 2.0, DECLARE, POWL, YAWL, and
Alignment-Based Conformance. Two working papers formalize mathematical
foundations: \texttt{runtime-verification-wasm.md} (Petri net soundness,
optimal alignment cost minimization, LTLf semantics for WASM execution) and
\texttt{declare-satisfaction-lattice.md} (5-valued dynamic valuation lattice
for streaming constraint checking). The \texttt{adversarial-canon.md} file
introduces Byzantine-fault-tolerant requirements that no existing canonical
process mining algorithm satisfies without modification.

\subsection{Primary Research Role}
\label{sec:sources-papers-identity-role}

The primary research role of this project within the dissertation is
\emph{academic-paper-classification-and-evidence-grounding}. It performs three
functions that no other component of the research foundry replicates. First, it
establishes the authoritative classification of 81 canonical process mining
and workflow papers across five mutually exclusive coverage statuses:
\texttt{COVERED\_BY\_TYPE}, \texttt{COVERED\_BY\_GRADUATION\_BOUNDARY},
\texttt{PARTIAL\_WITH\_REASON}, \texttt{DUPLICATE\_OR\_BACKGROUND}, and
\texttt{OUT\_OF\_SCOPE\_WITH\_REASON}. Second, it formalizes the Paper-to-Law
mapping algorithm $\mathcal{M}: I_P \to (\mathcal{T}_{\text{Rust}},
\mathcal{W})$, which converts formal structural invariants from published
literature into compile-time type-safety guarantees and non-forgeable witness
markers. Third, it documents the complete type-law obligation backlog --- 22
named missing types across 10 \texttt{PARTIAL} papers and 14 named missing
output shapes blocking algorithm graduation --- which drives the research
program's open work queue.

\subsection{Key Artifacts}
\label{sec:sources-papers-identity-artifacts}

The key artifacts manufactured by this project are:

\begin{itemize}
  \item \textbf{paper-canon.md} (v30.1.2, COMPLETE) --- the conformance
    registry for 9 core papers, each with formal objects, algorithms, failure
    conditions, compat type-law obligations, and board claim relevance.
  \item \textbf{PAPER\_TO\_TYPE\_LAW.md} --- 18 papers fully covered with Rust
    types; 39 at graduation boundary; 10 partial with named missing types.
  \item \textbf{PAPER\_TO\_EXECUTION\_LAW.md} --- 29 algorithms enumerated
    with input/output type shapes; all 29 currently carry \texttt{NO} in the
    \texttt{wasm4pm Coverage} column, documenting the graduation obligation.
  \item \textbf{PAPER\_TO\_BOARD\_CLAIM.md} --- 5+ Tier-1 structural
    completeness claims with named evidence paths into \texttt{wasm4pm-compat}
    source files and compile-fail fixtures.
  \item \textbf{PAPER\_CANON\_RECEIPT} --- cryptographic receipt certifying
    the 81-paper classification workflow, witnessed by the van der Aalst corpus
    and IEEE/ACM process mining bibliography, with zero gaps remaining.
  \item \textbf{declare-satisfaction-lattice.md} --- the 5-valued dynamic
    valuation lattice $\mathcal{V} = \{\bot, \text{PossiblySatisfied},
    \text{Satisfied}, \text{Violated}, \top\}$ with a verified join-semilattice
    structure for streaming DECLARE constraint evaluation.
  \item \textbf{core\_lifecycle\_algorithms.md} --- LaTeX-formatted mathematical
    foundations for four core lifecycle algorithms governing compilation,
    verification, admission, and cryptographic auditing.
\end{itemize}
